\section{Introduction}
Predicting case outcomes based on judge-summarized narratives is an important task. Unlike previous studies~\cite{malik-etal-2021-ildc, nigam-etal-2024-legal} and \cite{vats2023llms}, we aim to simulate realistic scenarios where legal judgment prediction systems are used to predict and explain judgments as cases arrive on the bench for adjudication. Our approach focuses on the core factual components of the case—specifically, the events that led to the case being filed, which serve as the basis for judgment prediction. These facts are the foundation of legal arguments and provide the context needed for making judicial decisions. In contrast to previous works that have included the entire case text (including proceedings), our focus on facts mirrors real-world conditions, where judges rely primarily on the case facts when delivering judgments.

In addition to the facts of the case, we incorporate additional legal information such as statutes, precedents, and arguments. Statutes represent codified legal principles, while precedents provide case-specific rulings that help guide decision-making. Together, these legal frameworks offer a structured basis upon which judges rely when formulating their rulings. By extracting and integrating these elements into our models, we aim to enhance both the prediction and explanation tasks by grounding the analysis in actual legal texts and the governing principles that are applied in real cases.

We explore the efficacy of various transformer-based models investigate the impact of summarizing legal judgments~\cite{xx4,xx9,xx10,xx11} using techniques~\cite{xx1,xx2,xx16,xx12} such as BERTSum~\cite{liu2019fine}, CaseSummarizer~\cite{polsley2016casesummarizer}, LetSum~\cite{farzindar2004atefeh}, and SummaRuNNer~\cite{nallapati2017summarunner}. Our findings suggest that leveraging summarized information yields decent results in judgment prediction. 

To further enhance the quality of prediction, we introduce hierarchical transformer models that utilize the entirety of judgment facts, demonstrating superior performance compared to traditional summarization methods. Additionally, our examination of LLMs, including Llama-2 (13b \& 70b) \cite{touvron2023llama} and GPT-3.5 Turbo \cite{brown2020language}, highlights the exceptional performance of GPT-3.5 Turbo in the context of Indian legal judgment prediction. We find that augmenting our models with additional legal information, such as statutes, precedents, and arguments, significantly improves the quality of both tasks.

In addition to focusing on the accuracy of legal judgment prediction, it is equally important to assess the quality of the explanations provided by the models. For this reason, we introduce two novel human evaluation metrics: \textit{Clarity} and \textit{Linking}. \textit{Clarity} refers to how well the predictions and explanations are structured and whether they convey the reasoning in a clear and understandable manner. This is critical in the legal domain, where complex legal concepts must be communicated effectively. \textit{Linking}, on the other hand, evaluates the logical consistency between the explanation and the final judgment. It assesses whether the explanation effectively ties back to the outcome and supports the predicted decision. These metrics are vital because, while models may produce accurate predictions, their explanations often lack coherence or fail to justify the decision meaningfully. By incorporating these metrics, we aim to ensure that models provide not only accurate outcomes but also transparent and interpretable explanations that can be trusted by legal professionals.


The key contributions of this study are:
\begin{enumerate}
    \item We focus on evaluating the performance of several transformer-based models and hierarchical transformer models, specifically on factual data, to mirror real-world conditions in judgment prediction. This approach contrasts with previous works that utilized full case texts.
    \item We utilize LLMs to assess their capabilities in legal judgment prediction and explanation tasks.
    \item We define two human evaluation metrics, \textit{Clarity} and \textit{Linking}, to assess the quality of LLM-generated judgment predictions and explanations, providing a comprehensive assessment of the overall task performance.
\end{enumerate}
To ensure reproducibility, both the code and dataset have been made publicly available via our repository\footnote{\href{https://github.com/ShubhamKumarNigam/Realistic_LJP}{https://github.com/ShubhamKumarNigam/Realistic\_LJP}}. Additionally, for convenience, we have uploaded the data\footnote{\href{https://huggingface.co/collections/L-NLProc/realistic-ljp-models-670ccbb2b390830b3989f6bf}{huggingface.co/collections/L-NLProc/Realistic\_LJP-models}} and models\footnote{\href{https://huggingface.co/collections/L-NLProc/realistic-ljp-datasets-670ccbeab5aea07a37e86df8}{huggingface.co/collections/L-NLProc/Realistic\_LJP-datasets}} to Huggingface.
